\documentclass[10pt,a4paper]{amsart}
\usepackage[margin=19mm]{geometry}
\usepackage{amsmath,amssymb,mathtools}
\usepackage[scr=rsfs,cal=euler]{mathalfa}
\usepackage{xfrac}
\usepackage[unicode,hidelinks,bookmarks=false]{hyperref}
\newcommand{\ie}{i.e.\ }
\newcommand{\cf}{cf.\ }
\newcommand{\resp}{respectively\ }
\newcommand{\Subsection}{\subsection}
\providecommand{\nobreakdash}{\nobreak\hbox{-}}
\setlength{\emergencystretch}{2em}
\multlinegap0pt
\usepackage[shortlabels]{enumitem}
\usepackage{bbm}
\usepackage{amssymb}
\usepackage{xcolor}
\usepackage{stackengine}
\usepackage[normalem]{ulem}
\usepackage{mathtools}
\newtheorem{lemma}{Lemma}
\newtheorem{theorem}{Theorem}
\newcommand\RR{\ensuremath{\mathbb{R}}}
\newcommand{\dd}{\ensuremath{\mathrm{d}}}
\title{Delay equations with measure-valued coefficients: From Carath\'eodory solutions to measurable semiflows}
\author{Marek Kryspin}
\date{Source: arXiv:2609.21098v1 (CC BY 4.0)}
\begin{document}
\maketitle

\noindent\emph{Source and licence.} Delay equations with measure-valued coefficients: From Carath\'eodory solutions to measurable semiflows. Version arXiv:2609.21098v1 (CC BY 4.0), distributed by arXiv under the Creative Commons Attribution 4.0 International licence, http://creativecommons.org/licenses/by/4.0/, as recorded in the arXiv licence metadata for this exact version and shown on the abstract page https://arxiv.org/abs/2609.21098v1. Full licence notice: \texttt{licenses/arxiv-2609.21098-CC-BY-4.0.txt}.\par\smallskip

\noindent\textit{Editorial scope.} Counted unit: the article's theorem thm:equivalent-parameter-measurability (source0.tex lines 2471--2499). The article's complete original proof of it is delivered with it (source0.tex lines 2501--2596, one \texttt{proof} environment of 96 lines). No mathematics is written, repaired or re-derived: the statement and the proof are the article's own lines, in the article's own notation, and no retained line is substituted or re-flowed.  Retained uncounted as the source-local context the proof needs: lines 2284--2305; lines 2407--2429.  Nothing was omitted from any retained span, so the package carries no omission record.  No retained line contains a citation, so the wrapper carries no bibliography and no bibliography entry.  No retained line contains a figure environment, an image-inclusion command or drawing code, so no image is delivered and the package declares no asset dependency.  Editorial formatting only, in addition: an AMS \texttt{amsart} preamble that restates the article's own declarations and macros used by the retained lines, namely the environments \texttt{lemma}, \texttt{theorem} on separate counters, together with the standard packages those lines need.  The retained environments print in this document as Lemma 1, Theorem 1 in order of appearance, whereas the article prints the same items with the numbering of its own full document.  The complete native source of the article as distributed in the arXiv e-print for this exact version is bundled unchanged as \texttt{sources/210936/source0.tex}.
\begin{lemma}\label{lem:measure-Borel}
For $z\in C(Q)$ and an open set $U\subset\mathbb R$,
define
\begin{equation*}
    B_{z,U}
    =
    \{\nu\in\mathcal M(Q):\langle z,\nu\rangle\in U\}.
\end{equation*}
Then
\begin{equation*}
    \mathfrak B(\mathcal M(Q),w^*)
    =
    \sigma\bigl(
        B_{z,U}:
        z\in C(Q),\
        U\subset\mathbb R\text{ open}
    \bigr).
\end{equation*}
Consequently, a mapping into $\mathcal M(Q)$ is
weak-$*$ Borel measurable if and only if all its pairings
with functions in $C(Q)$ are measurable.
\end{lemma}

\begin{lemma}\label{lem:delay-test-approximation}
For every $z\in C(Q)$, there exists a sequence
of functions
\begin{equation*}
    z_n(\tau,s)
    =
    \sum_{k=0}^{n-1}
    \mathbbm 1_{I_{k,n}}(\tau)\,x_{k,n}(\tau+s),
    \qquad
    x_{k,n}\in AC([-1,T]),
\end{equation*}
converging uniformly to $z$ on $Q$, where
\begin{equation*}
    I_{k,n}
    =
    \begin{cases}
        [kT/n,(k+1)T/n),
        & 0\le k<n-1,\\
        [(n-1)T/n,T],
        & k=n-1.
    \end{cases}
\end{equation*}
\end{lemma}

\begin{theorem}\label{thm:equivalent-parameter-measurability}
Let $\Sigma$ be a $\sigma$-algebra on $\mathfrak m_T$.
The following conditions are equivalent:
\begin{enumerate}[label=\textup{(}\roman*\textup{)}, ref=\textup{(}\emph{\roman*}\textup{)}]
\item\label{thm:equivalent-parameter-measurability-i}    For every $x\in AC([-1,T])$ and every
    $t\in[0,T]$, the mapping
    \begin{equation*}
        \mathfrak m_T\ni\mu
        \mapsto
        \int_0^t\int_{[-1,0]}
        x(\tau+s)\,\dd\mu_\tau(s)\,\dd\tau
        \in\mathbb R
    \end{equation*}
    is $(\Sigma,\mathfrak B(\mathbb R))$-measurable.
\item\label{thm:equivalent-parameter-measurability-ii}The mapping
    \begin{equation*}
        \mathfrak m_T\ni\mu
        \mapsto\nu(\mu)\in\mathcal M(Q)
    \end{equation*}
    is $(\Sigma,\mathfrak B(\mathcal M(Q),w^*))$-measurable.
\item\label{thm:equivalent-parameter-measurability-iii}The mapping $\mu\mapsto\nu(\mu)$ is weak-$*$ measurable:
    for every $z\in C(Q)$, the scalar mapping
    \begin{equation*}
        \mathfrak m_T\ni\mu
        \mapsto\langle z,\nu(\mu)\rangle\in\mathbb R
    \end{equation*}
    is $(\Sigma,\mathfrak B(\mathbb R))$-measurable.
\end{enumerate}
\end{theorem}

\begin{proof}
The equivalence of \ref{thm:equivalent-parameter-measurability-ii} and \ref{thm:equivalent-parameter-measurability-iii} follows directly
from Lemma~\ref{lem:measure-Borel}.
It remains to prove the equivalence of \ref{thm:equivalent-parameter-measurability-i} and \ref{thm:equivalent-parameter-measurability-iii}.

\medskip
\noindent
\ref{thm:equivalent-parameter-measurability-iii}$\Rightarrow$\ref{thm:equivalent-parameter-measurability-i}.
Fix $x\in AC([-1,T])$ and $t\in[0,T]$.
Choose continuous functions $\chi_n:[0,T]\to[0,1]$
by setting
\begin{equation*}
    \chi_n(\tau)
    =
    \begin{cases}
        1,&\tau\le t,\\
        \max\{1-n(\tau-t),0\},&\tau>t.
    \end{cases}
\end{equation*}
Then $(\tau,s)\mapsto\chi_n(\tau)x(\tau+s)$ belongs to $C(Q)$. Condition \ref{thm:equivalent-parameter-measurability-iii} implies that
\begin{equation*}
    \mathfrak{m}_T\ni \mu\mapsto
    \int_0^T\int_{[-1,0]}
    \chi_n(\tau)x(\tau+s)\,
    \dd\mu_\tau(s)\,\dd\tau \in \RR
\end{equation*}
is $(\Sigma,\mathfrak B(\mathbb R))$-measurable. For each fixed $\mu$, the inner integrals multiplied
by $\chi_n(\tau)$ converge to
\begin{equation*}
    \mathbbm 1_{[0,t]}(\tau)
    \int_{[-1,0]}x(\tau+s)\,\dd\mu_\tau(s).
\end{equation*}
Their absolute values are bounded by the integrable
function
\begin{equation*}
    \tau\mapsto \|x\|_{C([-1,T])}\|\mu_\tau\|_{\mathrm{TV}}. 
\end{equation*}
The dominated convergence theorem therefore gives
\begin{equation*}
    \int_0^T\int_{[-1,0]}
    \chi_n(\tau)x(\tau+s)\,
    \dd\mu_\tau(s)\,\dd\tau
    \to
    \int_0^t\int_{[-1,0]}
    x(\tau+s)\,\dd\mu_\tau(s)\,\dd\tau.
\end{equation*}
The limit is $(\Sigma,\mathfrak B(\mathbb R))$-measurable, proving \ref{thm:equivalent-parameter-measurability-i}.

\medskip
\noindent
\ref{thm:equivalent-parameter-measurability-i}$\Rightarrow$\ref{thm:equivalent-parameter-measurability-iii}.
Fix $z\in C(Q)$.
By Lemma~\ref{lem:delay-test-approximation}, there exist
functions
\begin{equation*}
    z_n(\tau,s)
    =
    \sum_{k=0}^{n-1}
    \mathbbm 1_{I_{k,n}}(\tau)x_{k,n}(\tau+s),
    \qquad x_{k,n}\in AC([-1,T]),
\end{equation*}
converging uniformly to $z$. For each $n$, consider
\begin{equation*}
     \mathfrak{m}_T\ni \mu\mapsto
    \int_0^T\int_{[-1,0]}
    z_n(\tau,s)\,\dd\mu_\tau(s)\,\dd\tau \in \RR.
\end{equation*}
This is a finite sum of mappings of the form
\begin{equation*}
     \mathfrak{m}_T\ni \mu\mapsto
    \int_a^b\int_{[-1,0]}
    x_{k,n}(\tau+s)\,\dd\mu_\tau(s)\,\dd\tau\in \RR.
\end{equation*}
Each summand is $(\Sigma,\mathfrak B(\RR))$-measurable by \ref{thm:equivalent-parameter-measurability-i}, since it
is the difference of two integrals with upper limits
$b$ and $a$. The choice of open or closed endpoints
does not affect these integrals, whose outer integration
is with respect to Lebesgue measure.

For every fixed $\mu\in\mathfrak m_T$, we have
\begin{equation*}
    \Big|
            \int_0^T\int_{[-1,0]}
            z_n(\tau,s)\,\dd\mu_\tau(s)\,\dd\tau
            -
            \langle z,\nu(\mu)\rangle
        \Big|
        \le
        \|z_n-z\|_{C(Q)}
        \int_0^T\|\mu_\tau\|_{\mathrm{TV}}\,\dd\tau
        \to0.
\end{equation*}
Thus $\mu\mapsto\langle z,\nu(\mu)\rangle$ is a pointwise
limit of $\Sigma$-measurable functions.
This proves \ref{thm:equivalent-parameter-measurability-iii} and completes the proof.
\end{proof}

\end{document}
